\documentclass[aspectratio=169,11pt]{beamer}
\usetheme{jcar}

%\includeonlyframes{a} % Para compilar solo en la que se está trabajando

%================================================
% PAQUETES
%================================================
\usepackage[utf8]{inputenc}
\usepackage[T1]{fontenc}
%\usepackage[spanish,es-nodecimaldot]{babel}
\usepackage{amsmath,amssymb}
\usepackage{booktabs}
\usepackage{tikz}

\title{Intervalos de confianza}
\author[Juan C. Arboleda Rivera]{Juan Camilo Arboleda Rivera}
\institute[U. de A.]{Instituto de Biología\\
  Facultad de Ciencias Exactas y Naturales\\
  Universidad de Antioquia\\
  juan.arboleda2@udea.edu.co}
\date{v: \today}

\begin{document}

%------------------------------------------------
\begin{frame}
\titlepage
\end{frame}

\begin{frame}
 \frametitle{Contenidos}
 \tableofcontents
\end{frame}

%------------------------------------------------
%\begin{frame}{Plan de trabajo: 120 minutos}
%\begin{center}
%\begin{tabular}{lc}
%\toprule
%\textbf{Tema} & \textbf{Duración}\\
%\midrule
%Motivación y conceptos fundamentales & 15 min\\
%Variabilidad y error estándar & 20 min\\
%Intervalos de confianza para una media & 25 min\\
%Pausa y preguntas & 10 min\\
%Ejemplos biológicos guiados & 20 min\\
%Proporciones y comparación de grupos & 15 min\\
%Ejercicios y evaluación final & 15 min\\
%\bottomrule
%\end{tabular}
%\end{center}
%\end{frame}

%================================================
% CONCEPTOS FUNDAMENTALES
%================================================
\section{Variabilidad y error estándar}

\begin{frame}{Una pregunta biológica}
Un equipo de ornitólogos desea conocer el peso promedio
de los tucanes de una reserva natural.

\begin{itemize}
    \item La reserva contiene cientos de tucanes.
    \item Pesar todos los individuos requiere mucho tiempo y esfuerzo de
    muestreo.
    \item Se seleccionan 15 tucanes para estudiarlos.
    \item La media observada es de 125 gramos.
\end{itemize}

\begin{block}{Pregunta}
¿La media de la muestra es exactamente igual a la
media de todos los tucanes de la reserva?
\end{block}

\end{frame}

%------------------------------------------------
\begin{frame}{Población, muestra y parámetro}
\begin{block}{Población}
Todos los organismos o unidades que nos interesa estudiar.
\end{block}

\begin{block}{Muestra}
Subconjunto de la población del que obtenemos datos.
\end{block}

\begin{block}{Parámetro}
Cantidad que describe una característica poblacional,
como la media $\mu$ o la proporción $p$.
\end{block}

\begin{block}{Estadístico}
Cantidad calculada con los datos de la muestra,
como la media $\bar{x}$ o la proporción $\hat{p}$.
\end{block}
\end{frame}

%------------------------------------------------
\begin{frame}{Ejemplos en Biología}
\begin{center}
\begin{tabular}{p{4cm}p{4cm}p{3cm}}
\toprule
\textbf{Pregunta} & \textbf{Parámetro} & \textbf{Estimador}\\
\midrule
Altura media de árboles & $\mu$ & $\bar{x}$\\
Proporción de semillas germinadas & $p$ & $\hat{p}$\\
Diferencia de crecimiento entre tratamientos
& $\mu_1-\mu_2$ & $\bar{x}_1-\bar{x}_2$\\
\bottomrule
\end{tabular}
\end{center}

\begin{alertblock}{Idea fundamental}
Con frecuencia conocemos la muestra, pero queremos
obtener información sobre una población más amplia.
\end{alertblock}
\end{frame}

%------------------------------------------------
\begin{frame}{¿Por qué existe incertidumbre?}
Imaginemos que tomamos varias muestras independientes
de una misma población.

\[
\bar{x}_1,\quad \bar{x}_2,\quad \bar{x}_3,\quad
\bar{x}_4,\quad \ldots
\]

Las medias muestrales generalmente serán diferentes.

\begin{itemize}
    \item Los organismos presentan variabilidad natural.
    \item Cada muestra contiene individuos diferentes.
    \item Una estimación puntual no describe por sí sola
    toda la incertidumbre.
\end{itemize}

Por eso necesitamos herramientas que acompañen la
estimación con una medida de precisión.
\end{frame}

%------------------------------------------------
\begin{frame}{Desviación estándar y error estándar}
La desviación estándar muestral describe la dispersión
de las observaciones individuales:

\[
s=\sqrt{\frac{\sum_{i=1}^{n}(x_i-\bar{x})^2}{n-1}}
\]

El error estándar de la media ($ SE $, de Standard Error) estima la
variabilidad de las medias muestrales:

\[
SE=\frac{s}{\sqrt{n}}
\]

Donde:
\begin{itemize}
    \item $s$: desviación estándar muestral.
    \item $n$: tamaño de la muestra.
\end{itemize}
\end{frame}

%------------------------------------------------
\begin{frame}{¿Qué ocurre cuando aumenta el tamaño muestral?}
Supongamos que la desviación estándar es de 20 gramos.

\begin{center}
\begin{tabular}{ccc}
\toprule
Tamaño muestral & Cálculo & Error estándar\\
\midrule
25 & $20/\sqrt{25}$ & 4 g\\
100 & $20/\sqrt{100}$ & 2 g\\
400 & $20/\sqrt{400}$ & 1 g\\
\bottomrule
\end{tabular}
\end{center}

\begin{block}{Conclusión}
Si la variabilidad se mantiene, aumentar el tamaño
de la muestra reduce el error estándar.
\end{block}

Una muestra mayor suele proporcionar estimaciones
más precisas, siempre que el diseño del estudio sea adecuado.
\end{frame}

%================================================
% INTERVALOS DE CONFIANZA
%================================================
\section{Intervalos de confianza}

\begin{frame}{¿Qué es un intervalo de confianza?}

Es un intervalo calculado a partir de una muestra
mediante un procedimiento estadístico que tiene una
determinada cobertura en muestreos repetidos.

\[
IC=\text{estimación}\pm\text{margen de error}
\]

El intervalo tiene dos límites:
\begin{itemize}
    \item Límite inferior.
    \item Límite superior.
\end{itemize}

\begin{block}{Ejemplo}
Un intervalo para la biomasa media podría ser
$[118,132]$ gramos.
\end{block}
\end{frame}

%------------------------------------------------
\begin{frame}{¿Qué significa un 95\% de confianza?}

\begin{block}{Interpretación}
Si repitiéramos el procedimiento de muestreo muchas
veces y calculáramos un intervalo en cada repetición,
aproximadamente el 95\% de los intervalos contendría
el verdadero parámetro poblacional.
\end{block}

\vspace{1em}

No significa que:
\begin{itemize}
    \item El 95\% de los datos esté dentro del intervalo.
    \item El parámetro fijo tenga necesariamente una
    probabilidad del 95\% de estar en el intervalo
    particular que ya calculamos.
\end{itemize}
\end{frame}

%------------------------------------------------
\begin{frame}
%[label=a]
{}{}
 \centering
 \includegraphics[height=.85\textheight]{confidence_interval.png}
\end{frame}

%------------------------------------------------
\begin{frame}
%[label=a]
{Visualización de intervalos}
\centering

\begin{tikzpicture}[x=0.7cm,y=0.55cm]
\draw[->] (0,0) -- (12,0) node[right] {Valor del parámetro};
\draw[dashed,red,thick] (6,-0.3) -- (6,6.5);
\foreach \y/\a/\b in {
1/3/8,
2/4/7,
3/2/5,
4/4/9,
5/3/6.8,
6/7/10
}{
\draw[thick,azul!50!black] (\a,\y)--(\b,\y);
\draw[thick,azul!50!black] (\a,\y-0.12)--(\a,\y+0.12);
\draw[thick,azul!50!black] (\b,\y-0.12)--(\b,\y+0.12);
}
\node[red,right] at (3.55,6.9) {\small Parámetro poblacional};
\end{tikzpicture}

\vspace{1em}
Cada línea representa un intervalo de una muestra.
Algunos intervalos pueden no contener el parámetro.
El nivel de confianza describe la cobertura a largo plazo.
\end{frame}

\begin{frame}
%[label=a]
{}{}
 \centering
 \includegraphics[width=\linewidth]{NYW-confidence-interval.pdf}\\
 La \textbf{media poblacional} ($ \mu $) estará dentro del intervalo de
 confianza calculado a partir de una muestra en el 95\% de los casos
\end{frame}

\begin{frame}
%[label=a]
{}{}
 \centering
 \includegraphics[height=\textheight]{interpreting-confidence-intervals-correctly.jpg}
\end{frame}

%------------------------------------------------
\begin{frame}
%[label=a]
{Intervalo de confianza del 95\% para una media}
Cuando la desviación estándar poblacional es desconocida,
o el tamaño de muestra es menor a 30, se utiliza habitualmente
la distribución $t$ de Student.

\[
IC_{95\%}=
\bar{x}\pm t_{0.975,n-1}\frac{s}{\sqrt{n}}
\]

Donde:
\begin{itemize}
    \item $\bar{x}$ es la media muestral.
    \item $s$ es la desviación estándar muestral.
    \item $n$ es el tamaño muestral.
    \item $t_{0.975,n-1}$ es el valor crítico de Student. También se usa
    $t_{0.025,n-1}$.
\end{itemize}

Para muestras pequeñas, es importante examinar si
los datos son aproximadamente normales y si existen
valores atípicos extremos.
\end{frame}

%----- Distribución t de Student --------------
\begin{frame}
%[label=a]
{}{}
 \centering
 \includegraphics[height=\textheight]{t-student-ic.png}
\end{frame}

\begin{frame}
%[label=a]
{}{}
 \centering
 \includegraphics[height=.95\textheight]{interval-of-confidence-t-dsitribution-graph.png}
\end{frame}

\section{Ejemplos}

%------------------------------------------------
\begin{frame}
%[label=a]
{Ejemplo: biomasa de plantas}
Se seleccionan 25 plantas y se obtiene:

\begin{itemize}
    \item Media: $\bar{x}=120$ g.
    \item Desviación estándar: $s=20$ g.
    \item Tamaño muestral: $n=25$.
    \item Valor crítico: $t_{0.975,24}=2.064$.
\end{itemize}

Primero calculamos el error estándar:

\[
SE=\frac{20}{\sqrt{25}}=4\text{ g}
\]

Después calculamos el margen de error:

\[
ME=2.064(4)=8.256\text{ g}
\]
\end{frame}

%------------------------------------------------
\begin{frame}{Resultado del intervalo}
\[
IC_{95\%}=120\pm8.256
\]

Por tanto:

\[
\boxed{IC_{95\%}=[111.74,\;128.26]\text{ g}}
\]

\begin{block}{Interpretación biológica}
El intervalo calculado para la media poblacional
de biomasa va aproximadamente de 111.74 a 128.26 g,
bajo los supuestos del método.
\end{block}

\begin{alertblock}{Pregunta}
¿Significa esto que el 95\% de las plantas pesa
entre 111.74 y 128.26 gramos?
\end{alertblock}

\pause
No. El intervalo estima la media poblacional, no
la distribución de los pesos individuales.
\end{frame}

%------------------------------------------------
\begin{frame}
%[label=a]
{Factores que afectan el intervalo}

La amplitud de un intervalo depende de:

\begin{enumerate}[<+->] \itemsep 0.5em
    \item \textbf{Variabilidad:} una desviación estándar
    mayor generalmente produce intervalos más anchos.
    \item \textbf{Tamaño muestral:} una muestra mayor
    generalmente reduce la amplitud.
    \item \textbf{Nivel de confianza:} un nivel más alto
    requiere normalmente un intervalo más ancho.
\end{enumerate}

\onslide<4->{
 \begin{block}{Pregunta de discusión}
 ¿Cómo podríamos obtener una estimación más precisa
 de la biomasa de una especie?
 \end{block}
}
\end{frame}

%------------------------------------------------
\begin{frame}
%[label=a]
{Ejemplo}{}
 \centering
 \includegraphics[height=.85\textheight]{smith-et-al-2016.png}
\end{frame}

\begin{frame}
%[label=a]
{}{}
 \centering \footnotesize
 \includegraphics[height=.8\textheight]{icad70149-fig-0002-m.jpg}\\
 Respuesta de (a) la riqueza de especies y (b) la abundancia total al uso de
 tierra e intensidad del uso de tierra. Las barras de error corresponden a
 intervalos de confianza del 95\%.
 \blfootnote{Smith, K. N. et al. Interactions of land use and climate change
 disproportionately impact Lepidoptera of the five major insect orders.
 Insect Conservation and Diversity n/a,.
 }
\end{frame}

%================================================
% PROPORCIONES
%================================================

%\begin{frame}{Intervalo de confianza para una proporción}
%Un equipo evalúa la germinación de 100 semillas.
%Germinan 70 semillas.
%
%La proporción observada es:
%
%\[
%\hat{p}=\frac{70}{100}=0.70
%\]
%
%Una aproximación normal al intervalo del 95\% es:
%
%\[
%\hat{p}\pm1.96
%\sqrt{\frac{\hat{p}(1-\hat{p})}{n}}
%\]
%
%Calculamos el error estándar:
%
%\[
%EE=\sqrt{\frac{0.70(0.30)}{100}}
%\approx0.0458
%\]
%
%Entonces:
%
%\[
%IC_{95\%}\approx0.70\pm1.96(0.0458)
%\]
%
%\[
%\boxed{IC_{95\%}\approx[0.610,0.790]}
%\]
%
%El intervalo corresponde aproximadamente a una
%proporción entre 61\% y 79\%.
%\end{frame}

%------------------------------------------------
%\begin{frame}{Comparación de dos tratamientos}
%Un experimento compara el crecimiento vegetal.
%
%\begin{center}
%\begin{tabular}{lcc}
%\toprule
% & Control & Fertilizante\\
%\midrule
%Media & 12 cm & 16 cm\\
%Desviación estándar & 4 cm & 5 cm\\
%Tamaño muestral & 25 & 25\\
%\bottomrule
%\end{tabular}
%\end{center}
%
%Diferencia entre las medias:
%
%\[
%\bar{x}_2-\bar{x}_1=16-12=4\text{ cm}
%\]
%
%Error estándar estimado de la diferencia:
%
%\[
%EE=\sqrt{\frac{s_1^2}{n_1}+\frac{s_2^2}{n_2}}
%\]
%
%\[
%EE=\sqrt{\frac{16}{25}+\frac{25}{25}}
%=\sqrt{1.64}\approx1.28
%\]
%\end{frame}
%
%%------------------------------------------------
%\begin{frame}{Intervalo para la diferencia de medias}
%Con una aproximación de Welch, el valor crítico
%para un intervalo del 95\% es aproximadamente 2.01.
%
%\[
%IC_{95\%}\approx4\pm2.01(1.28)
%\]
%
%\[
%\boxed{IC_{95\%}\approx[1.43,\;6.57]\text{ cm}}
%\]
%
%\begin{block}{Interpretación}
%Los datos son compatibles con un mayor crecimiento
%medio en el grupo fertilizado que en el grupo control.
%\end{block}
%
%La interpretación depende de los supuestos estadísticos
%y del diseño experimental. Un intervalo que no incluye
%cero sugiere una diferencia estadísticamente detectable,
%pero no demuestra por sí solo causalidad ni relevancia
%biológica.
%\end{frame}

%================================================
% ACTIVIDAD
%================================================

\section{Ejercicios}

\begin{frame}{Ejercicio 1}
Un grupo mide el diámetro de 25 árboles.

Datos:
\begin{itemize}
    \item Media muestral: 18 cm.
    \item Desviación estándar: 5 cm.
%    \item Valor crítico aproximado: $t^*=2.03$.
\end{itemize}

Resuelvan:
\begin{enumerate}
    \item Calculen el error estándar.
    \item Calculen el margen de error.
    \item Obtengan el intervalo del 95\%.
    \item Interpreten el resultado biológicamente.
    \item Expliquen qué ocurriría si la muestra tuviera
    144 árboles y la misma desviación estándar.
\end{enumerate}
\end{frame}

%------------------------------------------------
\begin{frame}{Solución de la actividad}
Primero calculamos el Error estándar:
\vspace{-.2em}
\[
SE=\frac{5}{\sqrt{25}}=1\text{ cm}
\]

Calculamos el Margen de error:
\vspace{-.2em}
\[
ME=2.064(1)=2.064\text{ cm}
\]

Intervalo:
\vspace{-.2em}
\[
IC_{95\%}=18\pm2.064
\]
\vspace{-.2em}
\[
\boxed{IC_{95\%}=[15.936,\;20.064]\text{ cm}}
\]

\begin{block}{Interpretación}
El intervalo estima la media poblacional del diámetro
entre aproximadamente 15.936 y 20.064 cm.
\end{block}
\end{frame}

%================================================
% EJERCICIO 2
%================================================

\begin{frame}{Ejercicio 2. Longitud de peces}
Un equipo de ecólogos estudia la longitud de una especie
de pez en un río de una región tropical.

Se selecciona una muestra aleatoria de 16 peces y se
obtienen los siguientes datos:

\begin{itemize}
    \item Longitud media muestral: $\bar{x}=12.4$ cm.
    \item Desviación estándar muestral: $s=2.0$ cm.
\end{itemize}

La desviación estándar poblacional es desconocida.

\medskip
\begin{alertblock}{}
    Calcule el intervalo de confianza del 95\% para
    la longitud media poblacional.
\end{alertblock}
\end{frame}

%------------------------------------------------
\begin{frame}{Ejercicio 2. Solución}
Como la desviación estándar poblacional es desconocida,
utilizamos la distribución $t$ de Student.
\vspace{-.2em}
\[
IC_{95\%}=\bar{x}\pm t_{\alpha/2,n-1}
\frac{s}{\sqrt{n}}
\]

Los grados de libertad son:
\vspace{-.2em}
\[
gl=n-1=16-1=15
\]

Para un nivel de confianza del 95\%:
\[
t_{0.025,15}\approx2.131
\]

Sustituyendo:
\[
IC_{95\%}=12.4\pm2.131\frac{2.0}{\sqrt{16}}
\]

Por tanto:
\[
IC_{95\%}=(12.4-1.0655,\;12.4+1.0655)
\]
\vspace{-1em}
\[
\boxed{IC_{95\%}=(11.33,\;13.47)\text{ cm}}
\]
\end{frame}

%----------------------
% EJERCICIO 3
%----------------------

\begin{frame}{Ejercicio 3. Longitud de insectos}
Un grupo de entomólogos estudia la longitud corporal de una
especie de escarabajo en un bosque tropical.

Se selecciona una muestra aleatoria de 12 escarabajos y se
obtienen los siguientes resultados:

\begin{itemize}
    \item Longitud media muestral: $\bar{x}=18.6$ mm.
    \item Desviación estándar muestral: $s=2.4$ mm.
\end{itemize}

\medskip
\textbf{Preguntas:}
\begin{enumerate}
    \item Calcule el intervalo de confianza del 95\% para
    la longitud media poblacional utilizando la distribución
    $t$ de Student.
    \item Determine los grados de libertad y el valor crítico
    correspondiente.
    \item Interprete el intervalo obtenido en el contexto
    del estudio entomológico.
\end{enumerate}
\end{frame}

\begin{frame}{Ejercicio 3. Solución: longitud de insectos}
Como la desviación estándar poblacional es desconocida,
utilizamos la distribución $t$ de Student.

La fórmula del intervalo de confianza es:

\[
IC_{95\%}=\bar{x}\pm t_{\alpha/2,n-1}
\frac{s}{\sqrt{n}}
\]

\textbf{Paso 1. Identificar los datos}

\[
n=12,\qquad \bar{x}=18.6,\qquad s=2.4
\]

\textbf{Paso 2. Calcular los grados de libertad}

\[
gl=n-1=12-1=11
\]
\end{frame}

\begin{frame}{Ejercicio 3. Solución: longitud de insectos}

Para un nivel de confianza del 95\%:
\[
\alpha=0.05,\qquad t_{0.025,11}=2.201
\]

\textbf{Paso 3. Calcular el margen de error}

\[
ME=2.201\frac{2.4}{\sqrt{12}}\approx1.525
\]

\end{frame}

\begin{frame}{Ejercicio 3. Solución e interpretación}
\textbf{Paso 4. Calcular el intervalo de confianza}

\[
IC_{95\%}=18.6\pm2.201\frac{2.4}{\sqrt{12}}
\]

\[
IC_{95\%}=18.6\pm1.525
\]

Por tanto:

\[
IC_{95\%}=(18.6-1.525,\;18.6+1.525)
\]

\[
\boxed{IC_{95\%}=[17.08,\;20.13]\text{ mm}}
\]
\end{frame}

\begin{frame}{Ejercicio 3. Solución e interpretación}
\medskip
\textbf{Interpretación biológica}\\~\\

Con un nivel de confianza del 95\%, estimamos que la
longitud media poblacional de esta especie de escarabajo
se encuentra entre 17.08 mm y 20.13 mm.

\end{frame}

%------------------------------------------------
\begin{frame}
%[label=a]
{Errores frecuentes}

\begin{itemize} \itemsep 1em
    \item Interpretar el intervalo como un rango que
    contiene al 95\% de los organismos.
    \item Pensar que una muestra grande elimina todos
    los posibles sesgos.
    \item Suponer que un intervalo estrecho implica
    que el resultado es biológicamente importante.
\end{itemize}

\end{frame}

%------------------------------------------------
\begin{frame}{Revisión de conceptos}

\begin{enumerate}[<+->] \itemsep 1em
    \item ¿Qué diferencia existe entre parámetro
    y estadístico?
    \item ¿Qué ocurre con el error estándar cuando
    aumenta el tamaño muestral?
    \item ¿Un intervalo estrecho garantiza que no
    existe sesgo?
    \item ¿Por qué es importante reportar intervalos
    de confianza en estudios biológicos?
\end{enumerate}
\end{frame}

%------------------------------------------------
\begin{frame}{Conclusiones}
\begin{itemize} \itemsep .5em
    \item La investigación biológica suele utilizar
    muestras para estudiar poblaciones.
    \item Toda estimación muestral tiene incertidumbre.
    \item El intervalo combina una estimación con
    un margen de error.
    \item El ancho depende de la variabilidad, el tamaño
    de la muestra y el nivel de confianza.
    \item La interpretación estadística debe acompañarse
    de una evaluación de la importancia biológica.
\end{itemize}

\vspace{0.5cm}
\begin{center}
\Large ¿Preguntas?
\end{center}
\end{frame}

%------------------------------------------------
\begin{frame}{Referencias}
\small
\begin{thebibliography}{9}

\bibitem{motulsky}
Motulsky, H. (2014).
\textit{Intuitive Biostatistics}.
Oxford University Press.

\bibitem{quinn}
Quinn, G. P. y Keough, M. J. (2002).
\textit{Experimental Design and Data Analysis for Biologists}.
Cambridge University Press.

\bibitem{zar}
Zar, J. H. (2010).
\textit{Biostatistical Analysis}.
Pearson.

\end{thebibliography}
\end{frame}

\end{document}
